\setcounter{section}{0}

\Needspace{5\baselineskip}
\hypertarget{prompt-based-methods}{%
\section{Prompt-Based Methods}\label{prompt-based-methods}}

Reasoning capabilities are becoming increasingly important in large models. Currently, model reasoning takes the form of continually generating tokens (tokens are also generated when a model displays that it is thinking, but they are wrapped in special markers and are therefore also billed in API calls), explicitly advancing logic through ``self-talk.'' The following are prompt-based methods for enhancing reasoning.

\Needspace{5\baselineskip}
\hypertarget{i.-chain-of-thought}{%
\subsection{I. Chain of Thought}\label{i.-chain-of-thought}}

Ask the model to think step by step, breaking a ``large problem'' into ``small steps.'' This shortens the distance each step must ``cross'' in embedding space, enabling complete reasoning by repeatedly ``predicting the next step.''

Chain of thought is the foundation of large-model reasoning techniques. Its limitations are that each generated step infers only ``forward'' from preceding content, without working backward from ``later objectives'' to ``earlier required steps''; it lacks self-correction and reflection mechanisms, allowing errors to accumulate; and it follows one path without exploration, making complex, diverse reasoning paths difficult to cover. These three limitations correspond to the three types of methods discussed later.

\Needspace{5\baselineskip}
\hypertarget{ii.-task-planning}{%
\subsection{II. Task Planning}\label{ii.-task-planning}}

\Needspace{4\baselineskip}
\begin{enumerate}
\def\labelenumi{\arabic{enumi}.}
\tightlist
\item
  Plan--execute paradigm
\end{enumerate}

Explicitly separate reasoning into ``planning'' and ``execution.''

Planning stage: rather than immediately calculating, the model first generates a structured plan from the original problem, listing the key steps required. This constructs an overall guiding framework.

Execution stage: the model uses each plan step as a separate input, generating chains of thought and results step by step.

Advantage: ``overview first, steps second'' reduces omissions and calculation errors in intermediate stages, significantly improving adaptation to unseen tasks, particularly in zero-shot settings.

\Needspace{4\baselineskip}
\begin{enumerate}
\def\labelenumi{\arabic{enumi}.}
\setcounter{enumi}{1}
\tightlist
\item
  Subtask decomposition
\end{enumerate}

For complex problems whose difficulty exceeds the examples, such as multistep mathematical reasoning, a least-to-most progression can be used. Decompose a complex problem into subproblems, arrange them in increasing difficulty, and solve them in sequence. When solving problem k, provide the answers to the first k-1 problems as context.

This establishes explicit dependencies across steps, using preceding results to support current reasoning and prevent logical breaks.

\Needspace{5\baselineskip}
\hypertarget{iii.-dynamic-adjustment-and-reflection}{%
\subsection{III. Dynamic Adjustment and Reflection}\label{iii.-dynamic-adjustment-and-reflection}}

If task planning lets the model ``look ahead,'' dynamic adjustment and reflection let it ``look back,'' reflecting on and correcting errors in past reasoning trajectories.

Initially, prompts implemented this behavior. Later, models learned directly through reinforcement learning to incorporate dynamic adjustment and reflection into their reasoning trajectories.

\FloatBarrier
\Needspace{5\baselineskip}
\hypertarget{references-119}{%
\subsection{References}\label{references-119}}

\vspace{0.25em}
\begin{itemize}
\tightlist
\item
  Wei, J., Wang, X., Schuurmans, D., et al.~(2022). \href{https://arxiv.org/abs/2201.11903}{Chain-of-Thought Prompting Elicits Reasoning in Large Language Models}. NeurIPS 2022.
\item
  Wang, L., Xu, W., Lan, Y., et al.~(2023). \href{https://aclanthology.org/2023.acl-long.147/}{Plan-and-Solve Prompting: Improving Zero-Shot Chain-of-Thought Reasoning by Large Language Models}. ACL 2023. DOI: 10.18653/v1/2023.acl-long.147.
\item
  Zhou, D., Schärli, N., Hou, L., et al.~(2023). \href{https://openreview.net/forum?id=WZH7099tgfM}{Least-to-Most Prompting Enables Complex Reasoning in Large Language Models}. ICLR 2023.
\item
  Shinn, N., Cassano, F., Gopinath, A., Narasimhan, K., \& Yao, S. (2023). \href{https://arxiv.org/abs/2303.11366}{Reflexion: Language Agents with Verbal Reinforcement Learning}. NeurIPS 2023.
\end{itemize}

\clearpage
